% Options for packages loaded elsewhere
\PassOptionsToPackage{unicode}{hyperref}
\PassOptionsToPackage{hyphens}{url}
%
\documentclass[
  11pt,
]{article}
\usepackage{amsmath,amssymb}
\usepackage{iftex}
\ifPDFTeX
  \usepackage[T1]{fontenc}
  \usepackage[utf8]{inputenc}
  \usepackage{textcomp} % provide euro and other symbols
\else % if luatex or xetex
  \usepackage{unicode-math} % this also loads fontspec
  \defaultfontfeatures{Scale=MatchLowercase}
  \defaultfontfeatures[\rmfamily]{Ligatures=TeX,Scale=1}
\fi
\usepackage{lmodern}
\ifPDFTeX\else
  % xetex/luatex font selection
\fi
% Use upquote if available, for straight quotes in verbatim environments
\IfFileExists{upquote.sty}{\usepackage{upquote}}{}
\IfFileExists{microtype.sty}{% use microtype if available
  \usepackage[]{microtype}
  \UseMicrotypeSet[protrusion]{basicmath} % disable protrusion for tt fonts
}{}
\makeatletter
\@ifundefined{KOMAClassName}{% if non-KOMA class
  \IfFileExists{parskip.sty}{%
    \usepackage{parskip}
  }{% else
    \setlength{\parindent}{0pt}
    \setlength{\parskip}{6pt plus 2pt minus 1pt}}
}{% if KOMA class
  \KOMAoptions{parskip=half}}
\makeatother
\usepackage{xcolor}
\usepackage[margin=25mm]{geometry}
\usepackage{color}
\usepackage{fancyvrb}
\newcommand{\VerbBar}{|}
\newcommand{\VERB}{\Verb[commandchars=\\\{\}]}
\DefineVerbatimEnvironment{Highlighting}{Verbatim}{commandchars=\\\{\}}
% Add ',fontsize=\small' for more characters per line
\newenvironment{Shaded}{}{}
\newcommand{\AlertTok}[1]{\textcolor[rgb]{1.00,0.00,0.00}{\textbf{#1}}}
\newcommand{\AnnotationTok}[1]{\textcolor[rgb]{0.38,0.63,0.69}{\textbf{\textit{#1}}}}
\newcommand{\AttributeTok}[1]{\textcolor[rgb]{0.49,0.56,0.16}{#1}}
\newcommand{\BaseNTok}[1]{\textcolor[rgb]{0.25,0.63,0.44}{#1}}
\newcommand{\BuiltInTok}[1]{\textcolor[rgb]{0.00,0.50,0.00}{#1}}
\newcommand{\CharTok}[1]{\textcolor[rgb]{0.25,0.44,0.63}{#1}}
\newcommand{\CommentTok}[1]{\textcolor[rgb]{0.38,0.63,0.69}{\textit{#1}}}
\newcommand{\CommentVarTok}[1]{\textcolor[rgb]{0.38,0.63,0.69}{\textbf{\textit{#1}}}}
\newcommand{\ConstantTok}[1]{\textcolor[rgb]{0.53,0.00,0.00}{#1}}
\newcommand{\ControlFlowTok}[1]{\textcolor[rgb]{0.00,0.44,0.13}{\textbf{#1}}}
\newcommand{\DataTypeTok}[1]{\textcolor[rgb]{0.56,0.13,0.00}{#1}}
\newcommand{\DecValTok}[1]{\textcolor[rgb]{0.25,0.63,0.44}{#1}}
\newcommand{\DocumentationTok}[1]{\textcolor[rgb]{0.73,0.13,0.13}{\textit{#1}}}
\newcommand{\ErrorTok}[1]{\textcolor[rgb]{1.00,0.00,0.00}{\textbf{#1}}}
\newcommand{\ExtensionTok}[1]{#1}
\newcommand{\FloatTok}[1]{\textcolor[rgb]{0.25,0.63,0.44}{#1}}
\newcommand{\FunctionTok}[1]{\textcolor[rgb]{0.02,0.16,0.49}{#1}}
\newcommand{\ImportTok}[1]{\textcolor[rgb]{0.00,0.50,0.00}{\textbf{#1}}}
\newcommand{\InformationTok}[1]{\textcolor[rgb]{0.38,0.63,0.69}{\textbf{\textit{#1}}}}
\newcommand{\KeywordTok}[1]{\textcolor[rgb]{0.00,0.44,0.13}{\textbf{#1}}}
\newcommand{\NormalTok}[1]{#1}
\newcommand{\OperatorTok}[1]{\textcolor[rgb]{0.40,0.40,0.40}{#1}}
\newcommand{\OtherTok}[1]{\textcolor[rgb]{0.00,0.44,0.13}{#1}}
\newcommand{\PreprocessorTok}[1]{\textcolor[rgb]{0.74,0.48,0.00}{#1}}
\newcommand{\RegionMarkerTok}[1]{#1}
\newcommand{\SpecialCharTok}[1]{\textcolor[rgb]{0.25,0.44,0.63}{#1}}
\newcommand{\SpecialStringTok}[1]{\textcolor[rgb]{0.73,0.40,0.53}{#1}}
\newcommand{\StringTok}[1]{\textcolor[rgb]{0.25,0.44,0.63}{#1}}
\newcommand{\VariableTok}[1]{\textcolor[rgb]{0.10,0.09,0.49}{#1}}
\newcommand{\VerbatimStringTok}[1]{\textcolor[rgb]{0.25,0.44,0.63}{#1}}
\newcommand{\WarningTok}[1]{\textcolor[rgb]{0.38,0.63,0.69}{\textbf{\textit{#1}}}}
\usepackage{longtable,booktabs,array}
\usepackage{calc} % for calculating minipage widths
% Correct order of tables after \paragraph or \subparagraph
\usepackage{etoolbox}
\makeatletter
\patchcmd\longtable{\par}{\if@noskipsec\mbox{}\fi\par}{}{}
\makeatother
% Allow footnotes in longtable head/foot
\IfFileExists{footnotehyper.sty}{\usepackage{footnotehyper}}{\usepackage{footnote}}
\makesavenoteenv{longtable}
\setlength{\emergencystretch}{3em} % prevent overfull lines
\providecommand{\tightlist}{%
  \setlength{\itemsep}{0pt}\setlength{\parskip}{0pt}}
\setcounter{secnumdepth}{5}
\ifLuaTeX
  \usepackage{selnolig}  % disable illegal ligatures
\fi
\usepackage{bookmark}
\IfFileExists{xurl.sty}{\usepackage{xurl}}{} % add URL line breaks if available
\urlstyle{same}
\hypersetup{
  pdftitle={Direct formal recognition, target bounds, and a genus-zero equation for P16.4},
  pdfauthor={Unverified research draft},
  hidelinks,
  pdfcreator={LaTeX via pandoc}}

\title{Direct formal recognition, target bounds, and a genus-zero
equation for P16.4}
\author{Unverified research draft}
\date{7 October 2026}

\begin{document}
\maketitle

\section{Scope and mathematical
status}\label{scope-and-mathematical-status}

The unchanged original P16.4 is \textbf{not solved} here. In particular,
this manuscript does not assert an unconditional all-genus curve
equation or an arbitrary-dimensional prescribed-target equation. It
proves a genus-zero equation, including nonsplit conics, relative to the
permitted split-projective-line, operator, descent and norming-reference
grants. It also supplies complete proofs of a cohomological base-change
lemma, direct FS recognition for an already actual positive formal pair,
finite graph interpolation with positive new vertices, explicit Cartier
calculations on actual charts, a residual-field identification lemma,
and general target estimates. It then identifies the geometric
applications and equation inputs not completed. These are unverified
drafts, not additional accepted grants.

The original domain retained throughout the equation discussion is a
complete real-valued non-Archimedean field \(K\) with characteristic and
residue characteristic zero, a smooth geometrically integral projective
\(K\)-variety \(X\), and an \textbf{actual} ample rational line bundle
\(L\). Put \(T=X^{\mathrm{an}}\), \(n=\dim X\), and \(V=(L^n)>0\). No
analytic point is removed. Raw target mass is \(\mu(T)=V\), without a
probability rescaling or field-degree multiplier. Poles are literal
values of actual metrics. Singular psh means precisely a decreasing
FS-net endpoint on the fixed bundle, excluding the identically
\(-\infty\) metric. Relative to a continuous psh reference \(\psi\),
write \(u\) for the scalar of \(\psi+u\). Its energy is

\[
E(u)=\inf_{v\in C(T)\cap\operatorname{PSH},\ v\geq u}E(v),
\qquad
E(v)=\frac1{n+1}\sum_{j=0}^n
 \int_Tv\,\operatorname{MA}(v^{[j]},0^{[n-j]}).
\]

Then \(\mathcal E^1=\{E>-\infty\}\) is an actual class of all-point
potentials, rather than a completion of graph functions.

Only the complete O217, D214, N220, T171, Q222 and C218 bodies and their
unchanged clarifications are caller grants. Their exact extracted bodies
and SHA256 bindings are preserved locally. O supplies the genuine mixed
all-field operator, all-test output integrability, extended decreasing
energy continuity before \(\mathcal E^1\), weighted arbitrary-net
continuity, and affine-segment calculus. C supplies strict-winner Borel
locality and comparison. D supplies invariant actual descent and the
continuous-test dual identity. N requires an actual continuous
single-atom reference which norms \textbf{every} FS potential. T
supplies the actual trivial-field equation. Q supplies the actual split
projective-line equation. Only O is needed for Theorems 7--8 and
Corollary 9. The genus-zero theorem uses O, D, N and Q in the exact
applications proved below, with T at the trivial boundary. Primary
publications and native235/native243/ROOT238--246 are mathematical data,
not extra grants.

\section{A base-change proof using a nonnegative flat Čech
complex}\label{a-base-change-proof-using-a-nonnegative-flat-ux10dech-complex}

\textbf{Lemma 1.} Let \(R\) be a Noetherian ring,
\(f:Z\to\operatorname{Spec}R\) a proper flat finitely presented
morphism, and \(M\) an invertible sheaf. If \(H^i(Z,M)=0\) for every
\(i>0\), then \(P=H^0(Z,M)\) is finite projective over \(R\). For every
\(R\)-algebra \(S\),

\[
 H^0(Z_S,M_S)=P\otimes_RS,
 \qquad H^i(Z_S,M_S)=0\quad(i>0).
\]

The same assertions hold after tensoring the Čech complex with any
\(R\)-module. No flatness of \(S\) is required.

\emph{Proof.} Properness implies quasi-compactness and separatedness.
Choose finitely many affine opens covering \(Z\). Every finite
intersection is affine, because \(Z\) is separated over the affine base.
The alternating Čech complex \(C^\bullet\), using strictly increasing
sets of indices, is bounded and concentrated in degrees
\(0,\ldots,r-1\). It computes sheaf cohomology of \(M\). On each affine
intersection \(U=\operatorname{Spec}A\), \(A\) is flat over \(R\), and
the module of sections of \(M|_U\) is projective of rank one over \(A\).
It is therefore \(R\)-flat. Finite products occurring in Čech terms are
finite direct sums and are flat.

The affine cover and all its intersections pull back under any base
change. Sections of a quasi-coherent sheaf on an affine scheme commute
with that base change. Consequently \(C^\bullet\otimes_RS\) is the Čech
complex of \(M_S\).

There is a completely elementary module proof of the tensor assertion.
At the top nonzero degree \(m\), put \(Z^m=C^m\), which is flat. Higher
cohomology vanishing gives short exact sequences
\(0\to Z^{q-1}\to C^{q-1}\to Z^q\to0\) descending from \(q=m\) to
\(q=1\). If the middle and right modules are flat, the kernel is flat,
and the sequence stays exact after tensoring with every module.
Induction shows that all \(Z^q\), including \(Z^0=P\), are flat, and
that tensoring the whole complex preserves exactness above degree zero
and its degree-zero kernel. If \(m=0\), this assertion is immediate.
Proper coherent cohomology over the Noetherian base makes \(P\) finite;
finite modules over a Noetherian ring are finitely presented. A finitely
presented flat module is projective. Taking \(Q=S\) and using the
pulled-back Čech computation proves the result. \(\square\)

The decisive property is \textbf{nonnegative flat Čech terms}. A perfect
complex with only degree-zero cohomology, without this property, need
not have flat degree-zero cohomology. For example
\([\mathbb Z\xrightarrow{2}\mathbb Z]\) in degrees \(-1,0\) represents
\(\mathbb Z/2\); tensoring with \(\mathbb Z/2\) introduces degree \(-1\)
cohomology. Thus higher vanishing alone would not justify arbitrary base
change.

The finiteness input here is an ordinary Noetherian proper-cohomology
theorem. In the polarized application below it follows directly from
Noetherian Serre theory, applied on the descended scheme. We do not
invoke the general proper-flat-perfect pushforward theorem to prove
Lemma 1.

\section{\texorpdfstring{Whole mod-\(a\) generation and direct
recognition}{Whole mod-a generation and direct recognition}}\label{whole-mod-a-generation-and-direct-recognition}

We state the exact domain of the recognition theorem, so that existence
of its geometric objects is not concealed.

Let \(K\) be complete, nontrivially real-valued, of characteristic and
residue characteristic zero, \(R=K^\circ\), and \(a\in R\) with
\(0<|a|<1\). Let \(\mathfrak X\) be an actual separated \(a\)-adic
formal scheme with a finite affine cover

\[
 \operatorname{Spf} B,\qquad B=R\langle T_1,\ldots,T_s\rangle/I,
\]

where \(I\) is finitely generated, the displayed quotient is
\(a\)-adically complete, and \(B\) is \(R\)-flat. Assume its
residue-field fibre is a proper reduced nodal curve. Assume its generic
analytic fibre is \textbf{already identified} with the analytification
of the original smooth geometrically integral projective \(K\)-curve
\(X\). Let \(\mathfrak M\) be an actual formal invertible sheaf whose
specified generic identification is the same original integral multiple
\(dL\), where \(d>0\) clears the rational bundle. Suppose
\(\deg(\mathfrak M|Y)>0\) for every irreducible residual component
\(Y\). No statement that every curve and bundle admits these objects is
built into this theorem.

\textbf{Theorem 2 (recognition on an actual formal pair).} For all
sufficiently large \(k\), finitely many formal global sections generate
\(\mathfrak M^k\). Their generic sections are algebraic global sections
of the same original \(kdL\). The metric of \(\mathfrak M/d\) is exactly
the zero-weight pure FS metric of those sections divided by \(kd\), at
\textbf{every actual analytic point}. The theorem has no genus
restriction, discreteness restriction, algebraic-closedness restriction
or residue-nilpotence assumption, within the explicitly stated
formal-pair domain.

\emph{Proof, reductions and ampleness.} Rank one makes the \(a\)-adic
and valuation topologies on \(R\) equivalent: the sequence \(|a|^j\) is
cofinal among positive radius bounds. Thus \(R=\varprojlim R/a^j\), and
the same complete topology is used for the charts. Restricted series
modulo \(a^j\) are polynomials, since only finitely many coefficients
survive. Conversely a restricted series all of whose coefficients lie in
\(a^jR\) is \(a^j\) times a restricted series. Consequently

\[
R\langle T\rangle/a^jR\langle T\rangle=(R/a^j)[T].
\]

The reductions \(X_j=\mathfrak X\bmod a^j\) are flat finitely presented
schemes over \(R_j=R/a^j\), by the finite generators of \(I\) and base
change of \(R\)-flatness. They are quasi-compact by the finite cover.

The spectrum of \(R_j\) has one prime. Indeed for every
\(b\in\mathfrak m_R\), rank one gives \(b^N\in a^jR\) for some \(N\).
Thus \(\mathfrak m_R/a^jR\) is a nil ideal. It need not be nilpotent
with a bounded exponent. Each residue-field inclusion
\(\operatorname{Spec}\widetilde K\hookrightarrow\operatorname{Spec}R_j\)
is an arbitrary thickening. The locally finitely presented morphism
\(X_j\to\operatorname{Spec}R_j\) has proper residue-field base change.
Properness across arbitrary thickenings, Stacks 37.3.4(8), makes \(X_j\)
proper. This theorem does not require a bounded nilpotence exponent.

Positive degree on every component makes the residual invertible sheaf
ample, by Stacks 33.44.15. Ampleness is invariant under arbitrary
thickenings, Stacks 37.3.2; hence \(M_1=\mathfrak M\bmod a\) is ample on
the \textbf{whole} \(X_1\). Alternatively fibre ampleness for the proper
morphism, Stacks 37.50.3, gives a neighbourhood of the sole base point,
which is the whole base.

\emph{Proof, descent and base change.} Express \(R_1\) as the filtered
union of its finitely generated \(\mathbb Z\)-subalgebras. The finitely
presented \(X_1\) and invertible \(M_1\) descend to some stage
\((Z_0,M_0)\) over a finite-type \(\mathbb Z\)-algebra
\(S_0\subset R_1\). Enlarge the stage so that properness, flatness and
relative ampleness hold there. The precise applications are Stacks
32.10.1 and 32.10.3 for the scheme, line bundle, inverse and evaluation
maps; 32.13.1 for properness; 32.8.7, using 10.168.1(3), for flat
finite-presentation descent; and 32.4.15 for ampleness. All transition
morphisms are affine, all bases are affine and
quasi-compact/quasi-separated, and the descended scheme remains finitely
presented. These are ordinary primary theorems in their stated domains;
they are not additional caller grants.

For clarity, the finite-data part of this descent has no new infinitary
assertion. Choose finite affine charts, finite equations and finitely
many principal localizations covering their overlaps. Line-bundle
trivializations and their inverse transition functions, the gluing maps,
their inverses, and their finitely many cocycle equalities are finite
data. Any equality true after the filtered colimit becomes true at a
later stage, by finite presentation. Inverse transition maps and
evaluation identities descend at the same later stage. This gives the
actual bundle, rather than a replacement positivity class. Preservation
of flatness and properness is the nontrivial part supplied by the exact
limit theorems just named.

\(S_0\) is Noetherian. Noetherian relative Serre generation and
vanishing give, for every sufficiently large \(k\), finite global
generators of \(M_0^k\), finite \(H^0\), and zero \(H^i\) for every
\(i>0\). A finite affine cover bounds the possible cohomological
degrees, so a single cutoff works. This is the use of Stacks 30.16.1 and
its projective Serre proof. It is applied to \(S_0\), \textbf{never to
\(R_1\)}. Lemma 1 now proves arbitrary base change of this vanishing to
\(R_1\). The pulled-back finite sections generate \(M_1^k\), and

\[
 H^1(X_1,M_1^k)=0.
\]

This is whole mod-\(a\) cohomology, not residue-fibre cohomology
substituted for it. The finite-projective \(H^0\) assertion follows from
Lemma 1, but the next lifting argument only needs this \(H^1\) vanishing
and the finite generators.

\emph{Proof, compatible lifting and complete local generation.} Put
\(M_j=\mathfrak M\bmod a^j\). \(R\)-flatness of the charts and local
freeness of \(\mathfrak M^k\) give a sheaf isomorphism

\[
 a^j\mathfrak M^k/a^{j+1}\mathfrak M^k\simeq M_1^k
\]

by multiplication by \(a^j\). Injectivity follows from
\(a\)-torsion-freeness before taking the quotient. There is therefore a
compatible exact sequence

\[
0\longrightarrow M_1^k\longrightarrow M_{j+1}^k
 \longrightarrow M_j^k\longrightarrow0.
\]

These sheaves can be viewed on the same underlying topological space.
\(H^1(X_1,M_1^k)=0\) makes every map on global sections
\(H^0(M_{j+1}^k)\to H^0(M_j^k)\) surjective. Starting with the finitely
many generators at level 1, choose compatible lifts at all levels.
Completeness of local frames identifies the formal sheaf with the
inverse limit of its reductions. Global sections commute with this
inverse limit, since they are an equalizer of the local restriction
maps. The compatible lifts consequently define actual formal global
sections \(s_1,\ldots,s_r\).

On a complete affine formal open where \(\mathfrak M^k\) has frame
\(e\), write \(s_i=f_ie\). Modulo \(a\), generation gives
\(1=\sum_i b_{i,0}f_i+a h_1\), for integral coefficients. Apply the same
mod-\(a\) decomposition to \(h_1\), then to its error, and continue. The
coefficients

\[
 b_i=\sum_{j\geq0}a^j b_{i,j}
\]

converge in the complete chart ring, and the remainder \(a^Nh_N\) tends
to zero: it belongs to \(a^N B\). Thus \(1=\sum_i b_i f_i\) exactly.
This is constructive complete division, without a Noetherian Nakayama
assertion on \(R/a\). It proves formal generation.

\emph{Proof, field comparison and exact metric.} The generic fibre and
bundle were already specified to be \(X^{\mathrm{an}}\) and \(kdL\).
Thus each lifted section is an analytic global section of that original
bundle. Projective \textbf{field} \(H^0\) GAGA turns it into an
algebraic global section of \(kdL\), with the specified identification
unchanged. One precise primary comparison route is Fujiwara--Kato,
literal arXiv:1308.4734v5, II.9.4.1--2. Its printed proof is in the
stronger topologically universally adhesive domain. The base here is in
that domain: \(R\) is \(a\)-adically complete and \(a\)-torsion-free,
\(K=R[1/a]\), and 0.9.2.7 and I.2.1.4(2) put the corresponding field
analytic setup in the proved domain. Every \(c\in K\) becomes integral
after multiplication by a sufficiently large power of \(a\), proving the
last field identity. Projectivity gives a proper finitely presented
algebraic scheme; the invertible sheaf is coherent; the comparison in
question is precisely
\(H^0(X,kdL)\to H^0(X^{\mathrm{an}},(kdL)^{\mathrm{an}})\). No general
carrier-algebraization theorem and no broader unproved-in-text
rigid-Noetherian comparison case is used. Conrad 2006 A.1.1 states the
same field comparison but refers its proof to Koepf 1974; that referral
is not treated as a printed proof here.

At every analytic point of this formal open, each integral \(f_i\) and
each integral \(b_i\) has absolute value at most one. The identity
\(1=\sum_i b_i f_i\) forces \(\max_i|f_i(x)|\geq1\); the reverse
inequality is already known. Thus \(\max_i|f_i(x)|=1\) at every point,
including supported points. In the formal frame \(\|e(x)\|=1\), so the
finite zero-weight FS system has exactly this metric. Generation also
rules out algebraic base points: a nonempty common zero locus after a
finite residue-field extension would have an analytic point,
contradicting the displayed identity. Divide by \(kd\) to get the metric
on the same original rational \(L\). \(\square\)

Theorem 2 proves a recognition theorem on an explicit domain of existing
formal pairs. It eliminates the extra hypotheses of whole-mod-\(a\)
vanishing, generation and section lifting in the direct route. It does
not manufacture the pair or its prescribed generic identification for
every curve. Its primary comparison and limit results remain ordinary
cited mathematical inputs; the source ledger separately records the
incomplete recursive audit of their foundations. Consequently no
universally applicable curve realization certificate is inferred from
this theorem.

\section{Horizontal and vertical Cartier calculations on actual
charts}\label{horizontal-and-vertical-cartier-calculations-on-actual-charts}

The following statements prove local and gluing implications, not the
existence of all the required charts on every original curve.

\textbf{Lemma 3 (a compatible horizontal section).} Let \(B\) be a
complete flat chart as above, let \(t\in B\) give a formally étale map
to the formal affine line near a smooth residual point, and let a given
bounded \(K\)-point \(p\) define evaluation \(B\to R\). Write
\(\tau=t-t(p)\). Suppose the residual zero fibre of \(\tau\), after
shrinking this actual chart, is just the reduced point \(\widetilde p\).
Then \((\tau)\subset B\) is an effective formal Cartier ideal,
\(B/(\tau)=R\), and its generic zero is exactly the given \(p\).

\emph{Proof.} For each \(j\), the zero fibre of \(\tau\) in
\(\operatorname{Spec}(B/a^j)\) is an étale finitely presented scheme
over \(R_j\), because it is the base change of the étale coordinate
chart along the section \(t=t(p)\). Its residue fibre is one reduced
rational point. Finiteness across arbitrary thickenings, Stacks
37.3.4(2), makes it finite. A finite étale algebra is finite locally
free. Over the local ring \(R_j\), the sole residue fibre has dimension
one, so it is free of rank one. Its unit element is a basis modulo the
maximal ideal and hence a basis; the structural map is an isomorphism to
\(R_j\). The given evaluation identifies that isomorphism and makes it
compatible for all \(j\).

Multiplication by \(t-t(p)\) is injective on \(R_j[t]\) for any ring
\(R_j\), by the monic polynomial leading-term calculation. Flatness of
the étale chart preserves that injectivity, so \(\tau\) is a
nonzerodivisor modulo every \(a^j\). It is a nonzerodivisor in \(B\) by
separation. If \(h\in B\) evaluates to zero, its image modulo every
\(a^j\) is uniquely \(\tau q_j\). Nonzerodivisor uniqueness makes the
\(q_j\) compatible. Completeness gives \(q\in B\) and \(h=\tau q\).
Therefore the kernel of evaluation is exactly \((\tau)\), and
\(B/(\tau)=R\). This also proves closedness of the ideal and
effectiveness of its inverse-limit Cartier presentation; no completion
of an unknown ideal is being assumed exact. The generic quotient is
\(K\), with the original evaluation, so its zero is precisely \(p\).
\(\square\)

On overlapping charts for this same section, two such generators have
the same principal ideal, and their ratio is a unit. Away from the
residual image of \(p\), the ideal is the unit ideal. They therefore
glue to an actual horizontal Cartier divisor. For finitely many
punctures reducing to distinct smooth points, the divisors are disjoint
and their signed integer sum is a formal Cartier divisor. If these
punctures are exactly the divisor of a meromorphic section \(\sigma\) of
the same original \(dL\), the generic bundle is identified with \(dL\)
by \(\sigma\). Each simple horizontal section meets one residual
component in one smooth rational point. Consequently its contribution to
residual degree is one on that component and zero on the others; signed
multiplicities give total residual degree \(d\deg L\).

To apply this universally, one still needs the actual adapted model, the
smooth relative étale coordinates and the bounded sections, with their
exact analytic/formal identification. BPR 3.13(3) and 4.8 give the
relevant puncture separation, while the smooth-coordinate/formal-fibre
inputs come from the original BL85. This manuscript has not completed
that entire recursive application. Genus zero starts with the split
projective line and its Gauss model, rather than a positive-genus
stable-model lemma.

\textbf{Lemma 4 (vertical gluing and degrees).} Work over an
algebraically closed nontrivially valued complete \(C\), with valuation
ring \(R\) and \(\Gamma=-\log|C^\times|\). Suppose an actual proper
strongly semistable formal model and a node/smooth-part cover of its
whole dual graph are given. Node \(e:v\to w\) has actual coordinates
\(xy=\pi_e\ne0\), length \(\ell_e=-\log|\pi_e|\in\Gamma\), and overlaps
where \(x\), respectively \(y\), is a formal unit. Let \(h\) have
\(\Gamma\)-valued vertex heights and integral slope \(k_e\) on every
edge, including both branches of every retained loop. Then an actual
vertical Cartier divisor with the specified generic trivialization
induces exactly \(h\), constant on complementary formal fibres. Its
degree on component \(Y_v\) is \(\sum_{e\ {\rm outgoing\ at} v}k_e\).

\emph{Proof.} Choose \(c_v\in C^\times\) with \(-\log|c_v|=h(v)\). Use
Cartier equations \(g_v=c_v\) on the smooth part and \(g_e=c_vx^{k_e}\)
on the node chart. These are elements of the total meromorphic ring,
nonzerodivisors, and generically invertible; negative exponents are
allowed in that ring. \(x\) and \(y\) are nonzerodivisors since
\(xy=\pi_e\) and \(R\)-flatness makes \(pi_e\) a nonzerodivisor. On the
first overlap \(g_e/g_v=x^{k_e}\). On the other,

\[
g_e/g_w=(c_v\pi_e^{k_e}/c_w)y^{-k_e}.
\]

The scalar has norm one because \(h(w)-h(v)=k_e\ell_e\); thus it and its
inverse are integral units. All overlap ratios are genuine formal units,
and the cocycle relations are literal ratios of the same meromorphic
equations. The local fractional invertible modules \(g_i^{-1}B_i\)
therefore glue. In their frame \(g_i^{-1}\), the canonical generic
section \(1\) has coefficient \(g_i\) and weight \(-\log|g_i|\), which
is \(h(v)+k_e(-\log|x|)\). Formal units have norm exactly one at every
generic analytic point because the unit and its inverse are both bounded
by one.

On the normalization of \(Y_v\), normalize the canonical section by
\(1/c_v\). Relative to the smooth-part frame it is nonvanishing. At a
node its coefficient is \(x^{k_e}\), with order \(k_e\) at that branch;
at the other branch the order is \(-k_e\). Summing orders gives the
stated degree. A self-loop contributes its two branches, whose sum is
zero. Subdivision leaves that cancellation unchanged. Rational slopes
are handled by first multiplying by a positive integer and then dividing
the resulting rational metric. Twisting an already specified extension
of \(dL\) by this divisor preserves that exact generic bundle through
the canonical trivialization. \(\square\)

This proof supplies the signed degree calculation. Nonnegative graph
coefficients alone still do not recognize a psh metric. Theorem 2
recognizes the exact twisted bundle when all its component degrees are
strictly positive.

\section{Rational interpolation preserving every cycle and positive
component}\label{rational-interpolation-preserving-every-cycle-and-positive-component}

\textbf{Theorem 5 (finite graph result).} Let \(G\) be a finite
connected metric graph, with loops permitted, whose edge lengths belong
to a nonzero divisible subgroup \(\Gamma\subset\mathbb R\). Let
\(\theta\) be any signed atomic vertex measure of total mass \(V>0\).
Let \(\mu\) be a positive atomic vertex measure of mass \(V\). A real PL
solution of \(\theta+\Delta F=\mu\) exists, where \(\Delta\) is the sum
of outgoing derivatives. It can be uniformly approximated, up to a
constant, by finite PL functions having shared \(\Gamma\)-vertex
heights, rational edge slopes, and one movable upward bend on each edge,
such that \textbf{every} vertex of their final graph has strictly
positive curvature for \(\theta+\Delta F_j\). Their total mass is \(V\).

\emph{Proof.} On the old vertices the weighted graph matrix has
quadratic form

\[
 -\sum_v F(v)\Delta F(v)=\sum_e\frac{(F(w_e)-F(v_e))^2}{\ell_e}.
\]

Connectivity implies that its kernel consists of constants; loops have
their correct zero endpoint contribution. Hence its image is the vectors
with sum zero. Solve the finite linear system for \(\mu-\theta\) and
interpolate linearly on each edge.

Let \(k\) be the number of \textbf{all} old vertices of this final
starting graph, after any preliminary loop subdivision or prescribed
type-II insertion. Put \(\rho=(V/k)\sum_v\delta_v\) and
\(\mu_\epsilon=(1-\epsilon)\mu+\epsilon\rho\). Its mass at every old
vertex is positive. Solve for \(F_\epsilon\), choosing one common
constant normalization; finite-dimensional linearity gives
\(F_\epsilon\to F\).

A nonzero divisible subgroup of \(\mathbb R\) is dense. Choose common
\(g_v\in\Gamma\) close to the vertex values of \(F_\epsilon\). On
oriented edge \(e\), let \(d=g_w-g_v\) and \(c=d/\ell_e\). Choose
rationals \(q_1<c<q_2\), both as close to \(c\) as desired. Put

\[
t_e=\frac{q_2\ell_e-d}{q_2-q_1}\in\Gamma,\qquad0<t_e<\ell_e.
\]

Use slope \(q_1\) from \(v\) to \(t_e\) and slope \(q_2\) from \(t_e\)
to \(w\). The endpoint equation \(q_1t_e+q_2(\ell_e-t_e)=d\) is exact.
Thus endpoints agree on every cycle, including irrational ones. At the
new bend, outgoing-slope curvature is \(q_2-q_1>0\). At each old vertex
the new outgoing slopes are arbitrarily close to those of
\(F_\epsilon\). There are finitely many of them, so choosing the errors
smaller than the minimum old positive margin divided by the maximum
valence keeps every old curvature positive. A loop counts both
incidences in that valence calculation. All new bends have strictly
positive mass.

The uniform difference from the chord of the chosen endpoints is at most
\(\ell_e(q_2-q_1)\), and the endpoint errors are uniformly small. This
proves uniform approximation. Edgewise opposite outgoing slopes cancel
in the total, so curvature mass remains \(V\). Take
\(\epsilon\downarrow0\) and the errors down sufficiently rapidly to
obtain the asserted approximating family. No vertex is inserted
afterwards; if a later actual construction requires another vertex,
restart the positivity step on that enlarged old graph. \(\square\)

Reference-only mixing would fail at a newly added zero-degree component.
The uniform old-vertex mixture is essential. Fixed-edge rational-slope
density also fails: on a circle with parallel edge lengths \(1\) and
\(\sqrt2\), rational linear slopes \(q,r\) with common endpoint height
must satisfy \(q=r\sqrt2\), forcing \(q=r=0\). The movable bends remove
this obstruction without discarding the loop.

For example take \(\theta=(1,1)\), \(\mu=(2,0)\), and \(\epsilon=1/10\).
The common endpoint height difference is \((9/10)(2-\sqrt2)\). Rational
slope pairs \((.52,.53)\) and \((.37,.38)\) on the edges of lengths
\(1,\sqrt2\) have bends \(-127+90\sqrt2\) and \(-180+128\sqrt2\). The
exact raw curvature masses are \(189/100,9/100,1/100,1/100\). Every mass
is positive and their sum is two. This is a finite graph calculation,
not evidence that an actual formal model or genuine measure has been
constructed.

The useful simplification of the direct route is that \(\theta\) in
Theorem 5 may be signed. An initial horizontal extension of the
meromorphic divisor of \(dL\) therefore need not reproduce a previously
given pure FS reference. Once its actual cover and generic
identification exist, Lemma 4 can construct the positive twists and
Theorem 2 can recognize them. This removes an unnecessary
reference-domination application, but does not replace the remaining
actual model construction or raw-curvature theorem.

\section{A residual-field bridge and the raw coefficient
boundary}\label{a-residual-field-bridge-and-the-raw-coefficient-boundary}

\textbf{Lemma 6 (residual field on a canonical chart).} Let \(C\) be as
in Lemma 4, \(B_C\) a reduced strict \(C\)-affinoid algebra with
spectral norm, and assume its canonical reduction is integral. Suppose
its Shilov boundary is the single actual point \(x\), and
\(|\mathcal H(x)^\times|=|C^\times|\). If an actual formal chart has
ring \(B_C^\circ\) and its special fibre is that canonical reduction,
then

\[
 \widetilde{\mathcal H(x)}=\operatorname{Frac}(\widetilde {B_C}).
\]

This is an equality through the natural reduction homomorphism, not a
conclusion from uniqueness of a reduction fibre alone.

\emph{Proof.} Single-Shilov norming gives \(\|f\|_{\rm sp}=|f(x)|\) for
every \(f\in B_C\). A nonzero canonical reduction of an integral-norm
element has value one at \(x\), so reduction embeds \(\widetilde B_C\)
into \(\widetilde{\mathcal H(x)}\), and hence embeds its fraction field.
The completed residue field is the completion of the fraction field of
\(B_C/\ker(x)\). A nonzero element \(f/g\) of norm one has
\(|f(x)|=|g(x)|\). By the value-group assumption choose
\(c\in C^\times\) with this common value. Both \(f/c\) and \(g/c\) then
have spectral norm one, so their reductions lie in \(\widetilde B_C\)
and the latter is nonzero. Its residue quotient is therefore in
\(\operatorname{Frac}(\widetilde B_C)\). Any norm-one element of
\(\mathcal H(x)\) has a fraction approximation with error strictly below
one and thus with the same residue. Norm-below-one elements have residue
zero. This proves surjectivity and the equality. \(\square\)

For the component point of a proper reduced semistable model, the
intended application takes an affine formal neighbourhood of its generic
component, removing all other components and nodes. The reduced special
fibre there is integral; its generic function field is \(\kappa(Y)\).
The canonical-chart recognition, single-Shilov identification and
type-II value-group assertion must first be established for that actual
chart. Once they are, Lemma 6 proves the required equality
\(\widetilde{\mathcal H(x_Y)}=\kappa(Y)\). BPR 4.2(3), 4.6 and 4.13
state the relevant canonical-model and component-point facts, but their
recursively cited reduction foundations have not all been reconstructed
in this audit. Lemma 6 cannot be applied just because a reduction fibre
has cardinality one.

For an actual formally induced metric already recognized in Theorem 2,
the correct primary coefficient chain is CLD \textbf{literal v1} 6.5.1,
6.6.8, and 6.9.3--6.9.4, with 6.9.5's calibrated-cone argument and
6.9.6's algebraically closed factor one. The germ reduction is the
Riemann--Zariski domain of the closure of the actual reduction. The
formal transition units identify the residual bundle with the pullback
of the restricted formal bundle. When the residual field is
\(\kappa(Y)\) and \(Y\) is a smooth proper residual curve, its proper
degree is the Riemann--Zariski degree. The resulting intended raw
formula is

\[
 \operatorname{MA}(\mathfrak M/d)
   =\sum_Y\frac{\deg(\mathfrak M|Y)}d\delta_{x_Y}.
\]

There is no division by \(V\), no field-degree factor, and no
dimension-one factorial. CLD 6.9.2 identifies atomic support; it does
\textbf{not} prove this coefficient formula by itself. BE 8.11's
citation to 6.9.2 must therefore be supplemented. The recognized pure FS
metric is globally approachable in the original projective carrier
through the permitted continuous theory; an unrecognized graph function
would not have that property. GH 1.2.13--14 defines the corresponding
actual formal-model degree measure and total mass. Its B.14--15 supplies
slope/retraction compatibility in its actual-model domain. These
statements are not identifications of an independently chosen graph
measure with an unknown genuine operator merely because their masses
agree.

This audit read the full coefficient proof units in CLDv1, including
their stable and algebraically closed branches. The calibrated-cone
foundation cited inside that proof, and the whole canonical-reduction
source chain needed for each generated component-point application, are
not recursively complete here. Accordingly the displayed formula is
recorded as the exact remaining application to discharge, not as a newly
certified universal geometric result. The signed horizontal and vertical
degree calculations above independently check its total-degree
normalization \textbf{when} it applies; they do not establish its
applicability.

\section{Intrinsic target uniformity in every dimension, relative only
to
O}\label{intrinsic-target-uniformity-in-every-dimension-relative-only-to-o}

Here no curve geometry or formal-pair hypothesis is used. Work in the
original projective domain, with O as the complete operator
prerequisite. Normalize actual potentials by \(\sup_Tu=0\), and put

\[
\begin{aligned}
B(u)&=-E(u)\geq0,\qquad A_\mu(u)=\int_T(-u)\,d\mu,\\
U_\mu&=\sup_{\substack{u\in\mathcal E^1,\ \sup u=0\\ B(u)\leq1}}A_\mu(u),\\
J(\mu)&=\sup_{v\in C(T)\cap\operatorname{PSH}}
 \left(E(v)-\int_Tv\,d\mu\right).
\end{aligned}
\]

\(\mu\) is a positive Radon measure of raw mass \(V\). Constants cancel
in \(J\). Suprema of actual nontrivial usc potentials are finite on the
compact carrier.

\textbf{Theorem 7 (uniformity and one direction of the dual criterion).}
Every actual \(\mathcal E^1\) potential is \(\mu\)-integrable if and
only if \(U_\mu<\infty\). \(J(\mu)<\infty\) implies both conditions. For
finite \(J\), the supremum over all actual \(\mathcal E^1\) tests equals
the continuous-test supremum \(J\). These statements hold in arbitrary
dimension and have no support restriction. They do not assert that
all-test integrability implies finite \(J\) without an additional proved
reference bound.

\emph{Proof, integrability implies uniformity.} Otherwise choose
normalized \(u_j\in\mathcal E^1\) with \(B(u_j)\leq1\) and
\(A_\mu(u_j)\geq4^j\). The actual convex combinations

\[
w_N=\sum_{j=1}^N2^{-j}u_j+2^{-N}0
\]

are nonpositive and decreasing. O's actual convexity and concavity give
\(E(w_N)\geq-1\). Actual decreasing closure gives a possibly identically
\(-\infty\) endpoint. The extended majorant-energy continuity argument
of O applies \textbf{before} membership in \(\mathcal E^1\) is known. An
identically \(-\infty\) endpoint would force the energies to tend to
\(-\infty\): for every finite \(M\), compactness puts a tail below
\(-M+1\), and monotonicity bounds its energy by \(V(-M+1)\). Thus the
endpoint \(w\) is actual nontrivial psh and \(E(w)\geq-1\). Sequential
monotone convergence, on the single measure \(\mu\), yields

\[
\int_T(-w)\,d\mu=\sum_{j\geq1}2^{-j}A_\mu(u_j)
 \geq\sum_{j\geq1}2^j=+\infty,
\]

a contradiction. The limit need not be sup-normalized or finite at a
preselected anchor; neither is used.

\emph{Proof, uniformity implies integrability.} Concavity gives
\(B(tu)\leq tB(u)\) for \(0<t\leq1\). The actual scaled potential is
still normalized and finite energy. Choose \(t=1\) when \(B(u)\leq1\),
otherwise \(t=1/B(u)\). Then \(tA_\mu(u)\leq U_\mu\), so \(A_\mu(u)\) is
finite. Constant shifts recover absolute integrability of every
unnormalized potential, since its positive part is bounded.

\emph{Proof, finite dual and actual tests.} Take continuous decreasing
\(v_i\downarrow u\in\mathcal E^1\). On a tail \(v_i\leq M\), a fixed
finite constant. Then the continuous nonnegative functions \(M-v_i\)
increase to \(M-u\). The dual inequality gives

\[
\int_T(M-v_i)\,d\mu\leq MV+J(\mu)-E(v_i)
 \leq MV+J(\mu)-E(u)<\infty.
\]

For completeness, the net integration fact used here is specific to
continuous, or nonnegative lower semicontinuous, functions on a compact
Hausdorff space with a Radon measure. At any strict level the
corresponding directed open sets cover the limiting open set; inner
approximation by a compact subset and a finite subcover give equality of
the limiting mass with the supremum of stage masses. One common index
approximates a finite lower layer sum at its finitely many levels.
Truncating height and refining those finite sums proves that the
integrals increase to the limiting integral. This is not arbitrary
measurable-net monotone convergence. The finite uniform bound therefore
proves integrability of \(u\). It also gives
\(\int v_i\,d\mu\to\int u\,d\mu\). O gives \(E(v_i)\downarrow E(u)\), so
the functional at \(u\) is a limit of continuous-test values and at most
\(J\). Continuous tests are themselves finite energy; the two suprema
agree. \(\square\)

\textbf{Theorem 8 (a sharper scaling estimate and conditional
coercivity).} For normalized actual \(u\in\mathcal E^1\), set

\[
A_k=\int_T(-u)\,\operatorname{MA}(u^{[k]},0^{[n-k]}),
\qquad D_n=\frac{n(n+1)}2.
\]

Then for all real \(0\leq t\leq1\),

\[
 B(tu)\leq tA_0+D_nt^2B(u).
\]

If the actual reference additionally satisfies

\[
 A_0=\int_T(-u)\,\operatorname{MA}(0)\leq M_0<\infty
\quad\text{for every normalized actual }u\in\mathcal E^1,
\]

then all-test integrability is equivalent to finite \(J\), and

\[
A_\mu(u)\leq U_\mu\left(1+2M_0+\sqrt{2D_nB(u)}\right),
\qquad
J(\mu)\leq U_\mu(1+2M_0)+\tfrac12D_nU_\mu^2.
\]

The variational functional tends to \(-\infty\) as normalized
\(B(u)\to+\infty\). Neither a bound \(M_0\) for every variety nor actual
variational attainment is proved here.

\emph{Proof.} O makes all \(A_k\) finite and nonnegative and gives
\(\sum_{k=0}^nA_k=(n+1)B(u)\). Its genuine affine curvature and segment
derivative laws give the finite polynomial identity

\[
B(tu)=\sum_{k=0}^n\binom nk A_k
 \int_0^t s^k(1-s)^{n-k}\,ds.
\]

At \(s=0\) only the reference term occurs; there is no zero coefficient
multiplied by a pole in the potential definition. Bound the \(k=0\) term
by \(tA_0\). For \(k\geq1\), \(A_k\leq(n+1)B(u)\), and

\[
\sum_{k=1}^n\binom nk s^k(1-s)^{n-k}
=1-(1-s)^n\leq ns.
\]

Integrating proves the estimate. It holds for all real \(t\) by the
actual real affine calculus in O, rather than by a finite numerical
check.

Under the additional reference bound, let
\(t=(1+2M_0+\sqrt{2D_nB(u)})^{-1}>0\). The linear term is at most
\(1/2\), and the quadratic term is at most \(1/2\). Thus \(B(tu)\leq1\),
and \(tA_\mu(u)\leq U_\mu\). This proves the integral bound. Subtract
\(B(u)\), maximize the resulting quadratic in \(\sqrt{B(u)}\), and use
Theorem 7 to obtain the displayed dual bound and coercivity. \(\square\)

If an \(M_0\) estimate is proved only for continuous normalized psh
potentials, it extends to actual finite-energy endpoints: for a
continuous decreasing net \(u_i\downarrow u\), compactness gives
\(\sup u_i\downarrow\sup u\). Apply the continuous estimate to
\(u_i-\sup u_i\); its integral differs from that of \(u_i\) by the
constant \(V\sup u_i\). The same continuous/lsc Radon-net fact passes to
the endpoint. This is a valid extension of a proved bound, not a proof
that such a bound exists.

For \(n=1\), the exact identity is \(B(tu)=tA_0+t^2(B(u)-A_0)\). When
\(A_0=0\), it is exactly quadratic, and the raw Green construction would
have energy \(-V^2I/2\), not \(-VI/2\).

\textbf{Corollary 9 (known outputs attain their own dual).} If
\(v\in\mathcal E^1\) and \(\mu=\operatorname{MA}(v)\), then
\(J(\mu)<\infty\), and the supremum over actual \(\mathcal E^1\) is
attained at \(v\). This requires no \(M_0\).

\emph{Proof.} O makes every test integrable against \(\mu\). For any
\(u\in\mathcal E^1\), concavity on the actual segment and its finite
derivative at \(v\) give

\[
E(u)\leq E(v)+\int_T(u-v)\,\operatorname{MA}(v).
\]

All integrals are finite on the common finite-value full-measure set
specified by O. Therefore the functional at \(u\) is at most its finite
value at \(v\). Decreasing continuous majorants of \(v\) converge both
in energy and in \(\mu\)-integral, by O's weighted law or the Radon-net
fact. Their objective values converge to the value at \(v\). Hence this
value is exactly \(J\). \(\square\)

This is a necessary target property and attainment for an
\textbf{already known output}. It does not solve an inverse equation for
an arbitrary intrinsic target, and it does not imply every-solution
all-point uniqueness.

\section{An unconditional genus-zero sector, including nonsplit
conics}\label{an-unconditional-genus-zero-sector-including-nonsplit-conics}

This section avoids the unfinished positive-genus formal geometry
completely. Its mathematical prerequisites are precisely the full grants
O, D, N and Q, with T for trivial valuation, and elementary algebraic
curve facts.

\textbf{Lemma 10 (maximum on finite atomic support on the projective
line).} Over a complete algebraically closed nontrivially valued \(C\),
let \(A\) be an actual ample rational bundle on \(\mathbb P^1_C\), of
degree \(V>0\). Let \(\psi\) be a genuine continuous psh reference with
curvature \(\lambda\) supported on a nonempty finite set \(S\) of
type-II or type-III points. Then every continuous psh scalar \(f\)
relative to \(\psi\) satisfies \(\sup f=\max_S f\). In particular this
holds for every actual FS scalar.

\emph{Proof.} Choose a coordinate Gauss root \(\zeta\) and the standard
degree-\(V\) reference. Q supplies the genuine continuous inverse atom
scalar \(a_y=-Vg_\zeta(\,\cdot\,,y)\) for every hyperbolic point \(y\),
with \(a_y(\zeta)=0\) and \(\operatorname{MA}(a_y)=V\delta_y\). The full
Q body proves the rooted-tree kernel, retraction, Lipschitz and
all-point assertions, for every rational degree; it is not just an
inverse-target criterion. Alternatively O's degree-one energy and
segment derivative give the identity

\[
 \int p\,\operatorname{MA}(q)-\int q\,\operatorname{MA}(p)
       =V(p(\zeta)-q(\zeta))
\]

for actual finite-energy scalars relative to that reference. To verify
it directly, compare the linear coefficient of the energy identity along
\((1-t)p+tq\) with the segment derivative at \(t=0\); all integrals
exist by O. Applying it with \(q=a_y\) gives, for a continuous psh
absolute metric of curvature \(\eta\), its relative scalar

\[
 p(y)=p(\zeta)-\int g_\zeta(x,y)\,d\eta(x).
\]

Apply this both to \(\psi\) and \(\psi+f\). The difference on all
hyperbolic points is

\[
 f(y)=f(\zeta)-\int g_\zeta(x,y)\,d(\eta-\lambda)(x),
 \qquad\eta=\operatorname{MA}(\psi+f)\geq0.
\]

Take the finite tree \(\Gamma\) spanning \(\zeta\) and \(S\). Q's
retraction identity identifies the restricted kernel with the kernel for
\(r_{\Gamma*}\eta\). It is continuous and uniformly 1-Lipschitz in its
tree variables. Thus \(f|_\Gamma\) is Lipschitz, and its elementary
graph distributional Laplacian is \(r_{\Gamma*}\eta-\lambda\): the
kernel has \(\Delta_y g_\zeta(x,y)=\delta_\zeta-\delta_{r_\Gamma x}\),
and both measures have mass \(V\). All signs here use outgoing
derivatives.

Put \(M=\max_Sf\) and \(w=(f-M)_+|_\Gamma\). Integration by parts on
finitely many intervals and at their vertices gives

\[
 \int_\Gamma w\,d(r_{\Gamma*}\eta-\lambda)
       =-\int_\Gamma |w'|^2\,dt.
\]

For justification with Lipschitz rather than smooth functions, the
interval derivative \(f'\) has bounded variation: its distributional
derivative is a finite signed measure, as follows directly by
integrating the kernel's one-jump derivatives. Approximate \(w\) on each
interval uniformly and with derivatives in \(L^1\), keeping its endpoint
values, by the usual convolution and endpoint adjustment. Stieltjes
integration by parts passes to this limit, since \(f'\) is bounded and
the second-derivative measure is finite. The Lipschitz positive-part
chain rule gives \(w'f'=|w'|^2\) almost everywhere. Summing intervals,
the outgoing endpoint terms are exactly the graph vertex terms, so no
boundary term is discarded. Since \(w=0\) on \(S\), its \(\lambda\)
integral vanishes and the left side is nonnegative. Hence \(w'=0\)
almost everywhere. Connectedness and \(w|_S=0\) imply \(w=0\) on
\(\Gamma\).

On an attached branch beyond \(\Gamma\), for a hyperbolic \(y\) the
kernel increment \(g_\zeta(x,y)-g_\zeta(x,r_\Gamma y)\) is nonnegative,
and is zero for every \(x\in S\). The displayed representation therefore
gives \(f(y)\leq f(r_\Gamma y)\). This proves the bound on every
hyperbolic point. Continuity and Q's actual full-carrier tree
description extend it to all terminal points, including type I and IV.
Thus the maximum is on \(S\). \(\square\)

The graph integration step uses only the finite rooted tree already
proved in Q, for already genuine continuous metrics. It neither treats
lisse as CLD-smooth nor defines a new operator by an inverse measure
construction.

\textbf{Theorem 11 (actual genus-zero equation).} Let \(K\) have the
unchanged original field domain, let \(X/K\) be a smooth geometrically
integral projective curve of geometric genus zero, and let \(L\) be any
actual ample rational line bundle, of raw degree \(V>0\). For a positive
Radon \(\mu\) on \(X^{\mathrm{an}}\), the following are equivalent:

\begin{enumerate}
\def\labelenumi{\arabic{enumi}.}
\tightlist
\item
  \(\mu(X^{\mathrm{an}})=V\) and every actual \(u\in\mathcal E^1(L)\) is
  absolutely \(\mu\)-integrable.
\item
  \(\mu(X^{\mathrm{an}})=V\) and the continuous-test dual
  \(J_\psi(\mu)\) is finite for any continuous psh reference \(\psi\).
\item
  \(\mu=\operatorname{MA}(\phi)\) for an actual
  \(\phi\in\mathcal E^1(L)\).
\end{enumerate}

Every solution is equal to every other at \textbf{every analytic point},
including their poles, modulo one finite real constant. The operator is
the genuine one in O, not a target-defined measure. The normalized
solution maximizes the actual energy-minus-integral functional. This
includes nonsplit conics, non-algebraically closed fields and
nondiscrete valuations. It is restricted to geometric genus zero.

\emph{Proof, an actual finite orbit.} For trivial valuation T already
gives the assertion with these exact smooth, projective, geometrically
integral and ample hypotheses. Suppose the valuation is nontrivial, and
put \(C=\widehat{\overline K}\) and
\(G=\operatorname{Gal}(\overline K/K)\). Choose a finite Galois
extension \(E/K\) over which \(X\) has a point. Such an extension exists
by taking the finite residue field of a closed point and its Galois
closure, since characteristic zero makes it separable.

Here is the elementary actual splitting argument. Field extensions
commute with cohomology by the affine Čech complex and flat tensoring,
so geometric genus zero gives \(H^1(X_E,\mathcal O)=0\), and geometric
integrality/properness gives \(H^0(X_E,\mathcal O)=E\). The latter can
also be checked after algebraic closure: a regular function on a proper
integral curve is constant, since its image in \(\mathbb P^1\) is closed
and misses infinity. For the smooth rational point \(p\), the Cartier
divisor has degree one and

\[
0\longrightarrow\mathcal O\longrightarrow\mathcal O(p)
 \longrightarrow\mathcal O(p)|_p\longrightarrow0.
\]

Its long exact sequence shows that \(H^0(\mathcal O(p))\) has dimension
two and surjects onto the one-dimensional fibre quotient. The canonical
section is nonvanishing away from \(p\), and a section with nonzero
fibre quotient is nonvanishing at \(p\). Thus this actual linear system
is basepoint-free. Its nonconstant map to \(\mathbb P^1_E\) pulls back
\(\mathcal O(1)\) to degree-one \(\mathcal O(p)\); it is finite of
degree one and is an isomorphism, because a finite birational map to the
normal curve \(\mathbb P^1_E\) is an isomorphism. This proves the needed
splitting without importing the full Riemann--Roch theorem. Stacks
53.10.4 is a primary scope cross-check, rather than an extra NA grant.

Choose this actual isomorphism over \(E\). After clearing rational
denominators, choose the actual line-bundle identification
\(L_E\simeq V\mathcal O(1)\), in
\(\operatorname{Pic}(X_E)\otimes\mathbb Q\), using the projective-line
Picard calculation in Q. These identifications can be made over a finite
Galois extension after increasing \(E\) if necessary; no unknown
analytic carrier is introduced.

Pull the standard degree-\(V\) FS metric to \(X_C\). Its genuine measure
is \(V\delta_y\) at the transported coordinate Gauss point, by Q. The
subgroup \(\operatorname{Gal}(\overline K/E)\) fixes that metric and
that actual point: it fixes the coordinates and acts isometrically on
coefficients, so preserves the Gauss seminorm. Thus the metric and point
have finite actual \(G\)-orbits. This gives a finite orbit without a
fixed point on a positive-genus cyclic skeleton, and without a general
pure-FS support theorem.

Choose an actual pure FS metric \(\theta_K\) on \(L\), using a globally
generated integral multiple; its pullback \(\theta_C\) is invariant.
Normalize all orbit metrics by \(\sup(\phi_y-\theta_C)=0\). Naturality
of the genuine continuous operator under the semilinear isometries
follows from its actual continuous construction in D/O; O passes it to
decreasing limits. Q's every-solution uniqueness shows that metrics
belonging to the same orbit point are identical after this
normalization. The normalized finite family is permuted by \(G\). Its
finite arithmetic average \(\psi_C\) is a genuine continuous psh metric,
is invariant, and in dimension one O's affine curvature law gives

\[
 \operatorname{MA}(\psi_C)=
      \frac V{|S|}\sum_{y\in S}\delta_y,
 \qquad S=G y.
\]

Only finite \textbf{continuous} averaging is used. D's actual continuous
descent gives a metric \(\psi_K\) on the same original \(L\); D/O
pushforward gives \(\operatorname{MA}(\psi_K)=V\delta_\xi\), where
\(\xi\) is the common image of \(S\).

\emph{Proof, every FS maximum and the intrinsic condition.} For any
original FS potential relative to \(\psi_K\), its pullback is a
continuous psh scalar relative to \(\psi_C\) and is \(G\)-invariant.
Lemma 10 says its maximum is attained on \(S\); invariance makes its
values on \(S\) all equal. D's surjective quotient therefore gives
\(\sup_T f=f(\xi)\). This proves \textbf{every} FS norming equality
required by N on the actual downstairs bundle and reference, as well as
the single-atom measure identity. N now applies with no additional
geometric premise and proves the equivalence of conditions 1 and 2,
together with coercivity. Bounded continuous changes of reference
preserve the absolute energy class and dual finiteness, by the exact
energy cocycle in O.

\emph{Proof, prescribed target and descent.} Let condition 2 hold. D's
unique invariant Radon lift \(\widetilde\mu\) has raw mass \(V\) and the
identical finite continuous-test dual upstairs relative to \(\psi_C\).
The chosen actual isomorphism \(X_C=\mathbb P^1_C\) and line-bundle
identification meet Q's complete rational-bundle domain; its theorem is
reference-independent. Q therefore supplies an actual upstairs
finite-energy solution for \(\widetilde\mu\) and uniqueness of every
solution at every point. Normalize its scalar relative to invariant
\(\psi_C\) by supremum zero. Each \(g\) transforms it to a solution of
the same invariant target with the same normalization. Q's all-point
uniqueness forces invariance. D descends this actual invariant singular
metric, preserving energy and all pole values. O pushes its genuine
measure to \(\mu\). This is singular \textbf{descent}, not singular
noninvariant averaging.

Conversely every actual output has the raw mass and all-test
integrability by O. For uniqueness, pull back any two original
solutions. Their upstairs genuine outputs are invariant by equivariance
and have pushforward \(\mu\) by O. D's invariant-lift uniqueness
identifies both with \(\widetilde\mu\). Q makes the two metrics differ
by one finite constant at every upstairs point. Surjectivity of the
actual quotient gives that equality at every original point. Corollary 9
identifies the finite maximizing variational value for each solution.
\(\square\)

The construction replaces variational compactness and envelope/outer
differentiation in this sector by Q's actual tree equation and D's
invariant descent. O supplies energy differentiation on actual affine
segments, C retains comparison in the same domain, and the proof above
establishes all-point uniqueness through Q and the surjective quotient.
It asserts no arbitrary-envelope theorem over the remaining fields.

\section{Exact remaining equation
obligations}\label{exact-remaining-equation-obligations}

The actual recognition theorem is useful progress on the direct route,
but R1--R5 are not all discharged.

\begin{longtable}[]{@{}
  >{\raggedright\arraybackslash}p{(\columnwidth - 2\tabcolsep) * \real{0.5000}}
  >{\raggedright\arraybackslash}p{(\columnwidth - 2\tabcolsep) * \real{0.5000}}@{}}
\toprule\noalign{}
\begin{minipage}[b]{\linewidth}\raggedright
Obligation
\end{minipage} & \begin{minipage}[b]{\linewidth}\raggedright
Result here and remaining application
\end{minipage} \\
\midrule\noalign{}
\endhead
\bottomrule\noalign{}
\endlastfoot
R1: actual model on the original curve & Primary BL85 stable reduction
and BPR model/skeleton units read in the recorded ranges. Their full
recursively generated formal-fibre, finite-map, genus and
affinoid-subdivision foundations are not completed. No new unconditional
semistable existence theorem is claimed. \\
R2: cofinal retractions and full curve topology & BPR 3.13 and 4.11
identify actual refinements, including punctures and loops, in their
domain. The whole topology/formal-fibre application for every original
curve and every point is not reconstructed. \\
R3: exact same-bundle recognition and curvature & Lemmas 1, 3, 4 and
Theorem 2 prove cohomology, lifting, and chart/gluing implications.
Theorem 5 fixes strict positivity. The exact actual global pair and the
raw component-degree coefficient application still have to be verified
for every approximation. \\
R4: kernels and every-point uniqueness & No abstract kernel or graph
completion is asserted actual PSH here. Kernels must first be recognized
as actual continuous psh differences. FS section norms and formal frames
must recover terminal type-I/type-IV and pole values; actual annulus
convexity must recover interior type-III values. Lisse is not
CLD-smooth. \\
R5: descent anchor & Finite positive support of an invariant actual pure
FS measure would give finite Galois orbits. Actual continuous inverse
atoms and every-solution uniqueness must precede finite continuous orbit
averaging and D descent. Every FS maximum on the resulting support/fibre
must then be proved before using N. None of these unproved antecedents
is granted by D or N. \\
\end{longtable}

The higher-dimensional target route also remains geometric. A finite
reference Shilov norming \textbf{set} does not bound each individual
normalized value below. The native243 common coefficient-Gauss
generic-incidence construction must have its initial residual-coordinate
hypotheses verified for every component of an actual ample same-bundle
embedding. Original section values, rather than whole residual fields,
need preservation. The auxiliary transcendental field step cannot use GM
Appendix A.7, whose domain is algebraic-completion extensions. An actual
full-core voltage bound on the restricted curve could yield \(M_0\); no
such uniform original-variety application is proved here.

The original replacements have the following precise status. The mixed
operator, all-test output integrability, energy law, decreasing weighted
nets and \textbf{affine-segment} differentiation are supplied by O.
Strict locality and comparison are supplied by C. Actual
normalized-potential compactness over the remaining general nontrivial
fields, or a complete substitute for it, is absent. General envelope
regularity/orthogonality and both signs of differentiation in outer
continuous obstacle directions are absent where a variational proof
would need them. Theorem 8 is coercivity only with an actually proved
\(M_0\), not compactness or attainment. No prescribed-target attainment
for every intrinsic target in dimension at least two, and no
every-solution equality at every analytic point there, has been proved.
T and Q retain their exact established boundary domains and do not
remove these gaps.

If genuine continuous all-genus atoms were completed, the proposed
arbitrary-target curve construction would use directed finite full-core
retractions, not one skeleton. With \(\mu=V\nu\),
\(u_G=-V\int j_G\,d\nu\), its raw constants must be
\(E(u_G)=-V^2I_G/2\), \(\int u_G\,d\mu=-V^2I_G\), and objective
\(V^2I_G/2\). Finite dual bounds would bound energy costs before passing
to \(\mathcal E^1\); only after O's extended continuity establishes
actual finite energy could its operator identify the target limit. This
paragraph records the necessary normalization and logical order, not a
proved construction. No arbitrary singular noninvariant average is used.

\section{Primary audit and
reproduction}\label{primary-audit-and-reproduction}

The exact 22 supplied frozen files and all six full grant bodies match
their caller bindings. The grants and clarifications are preserved
byte-bound locally; all other source bytes remain primary data.
\texttt{source-application-ledger.md} lists the whole units read,
applications, and recursive audit gaps. The BL85 scan check actually
inspected printed pages 374--377 (PDF pages 27--30). It verifies the
separate stable \(g\geq2\) and semistable \(g\leq1\) clauses, the
\(g\geq1\) scope of Lemmas 7.2--7.3, the auxiliary integer prime to
residue characteristic, and the explicit omitted
affinoid-\(\mathbf P^1\) subdivision referral. Machine OCR is not
promoted to a checked formula on other uninspected scans.

The small standard-library program \texttt{finite\_checks.py} uses exact
rational arithmetic and an exact ordered \(\mathbb Q(\sqrt2)\)
implementation. It checks the stated irrational-cycle interpolation, a
specified finite rational slope grid, the Bernstein scaling coefficients
in dimensions 1--8 at \(t=j/20\), finite Čech monomial calculations over
four nonreduced rings, nilradical witnesses, and raw curve constants.
Exhaustiveness applies only to those explicit finite grids. It proves no
actual formal scheme, analytic comparison, psh theorem or inverse
equation.

Reproduce in this directory:

\begin{Shaded}
\begin{Highlighting}[]
\ExtensionTok{python3}\NormalTok{ verify\_packet.py}
\ExtensionTok{python3}\NormalTok{ finite\_checks.py}
\ExtensionTok{python3}\NormalTok{ build\_deliverable.py}
\end{Highlighting}
\end{Shaded}

\texttt{build\_deliverable.py} records commands, return codes, logs,
environment versions and artifact hashes, creates the standalone
\texttt{P16.4-paper.tex} from this full argument with existing pandoc,
and compiles twice using existing pdfLaTeX with shell escape disabled.
No installation or resource-heavy computation is needed. No supervised
compute facility was exposed, and none was launched. Primary PDF files
remain at the exact caller-authorized read-only packet paths; all
generated work is in this artifact directory. No account,
authentication, runtime, proxy, CA, sandbox setting, protected editor,
other campaign stream or GitHub mutation was changed or accessed.
Execution and compilation success are observations, not mathematical
correctness certificates. \texttt{candidate=false} remains required for
the unchanged original task.

\section{Primary identities}\label{primary-identities}

The six byte-bound caller bodies O217/D214/N220/T171/Q222/C218,
including their full scope limitations and harmless clarifications, are
the only accepted mathematical grants.

Bosch--Lütkebohmert, \emph{Stable reduction and uniformization of
abelian varieties I}, Math. Ann. 270 (1985), 349--379, original scanned
primary:
\url{https://gdz.sub.uni-goettingen.de/download/pdf/PPN235181684_0270/LOG_0034.pdf}.

Baker--Payne--Rabinoff, \emph{On the structure of nonarchimedean
analytic curves}, literal arXiv:1404.0279v1:
\url{https://arxiv.org/pdf/1404.0279v1}.

Fujiwara--Kato, \emph{Foundations of Rigid Geometry I}, literal
arXiv:1308.4734v5: \url{https://arxiv.org/pdf/1308.4734v5}. This is
distinct from Freitag--Kiehl, \emph{Étale cohomology and the Weil
conjecture} (1988), the item called {[}FK88{]} in GJKM.

Chambert-Loir--Ducros, \emph{Formes différentielles réelles et courants
sur les espaces de Berkovich}, literal arXiv:1204.6277v1:
\url{https://arxiv.org/pdf/1204.6277v1}.

Gubler--Hertel, \emph{Local heights of toric varieties over
non-Archimedean fields}, literal arXiv:1512.06574v2:
\url{https://arxiv.org/pdf/1512.06574v2}.

Boucksom--Eriksson, \emph{Spaces of norms, determinant of cohomology and
Fekete points in non-Archimedean geometry}, literal arXiv:1805.01016v3:
\url{https://arxiv.org/pdf/1805.01016v3}.

Gubler--Martin, \emph{On Zhang's semipositive metrics}, Documenta Math.
24 (2019), 331--372:
\url{https://ems.press/content/serial-article-files/26466}.

Conrad, \emph{Relative ampleness in rigid geometry}, AIF 56 (2006),
1049--1126: \url{https://www.numdam.org/item/10.5802/aif.2207.pdf}. Its
Koepf referral is recorded as such.

Stacks Project primary applications are linked by stable tag in the
source ledger. No summaries or historical review flags replace their
exact hypotheses or a missing proof.

\end{document}
